\documentclass[11pt]{article}
\usepackage{preamble}

\newcommand{\IconCheck}{[CHECK]}
\newcommand{\IconMic}{[MIC]}
\newcommand{\IconClipboard}{[CLIPBOARD]}
\newcommand{\IconChart}{[CHART]}
\newcommand{\lessonrow}[2]{\textbf{#1} & #2 \\ \hline}

\newcommand{\phasetag}[1]{%
  {\small Tag: #1\par}
}

\newcommand{\phaseitems}[1]{%
  \begin{itemize}[leftmargin=1.5em,itemsep=2pt,topsep=2pt]
  #1
  \end{itemize}%
}

\newcommand{\teacheranswerbox}[2]{%
  \begin{callout}{#1}{calloutgreen}
  #2
  \end{callout}
}

\newcommand{\visualplaceholder}[1]{%
\begin{center}
\begin{tcolorbox}[
  enhanced,
  width=0.9\linewidth,
  colback=frameshade,
  colframe=framegray!45,
  boxrule=0.5pt,
  arc=2pt,
  left=6pt,
  right=6pt,
  top=5pt,
  bottom=5pt
]
\small #1
\end{tcolorbox}
\end{center}
}

\newcommand{\entrybox}{\framebox[1.1em]{\rule{0pt}{0.95em}}}

\newcommand{\qfortyfourtable}{%
\begin{center}
\renewcommand{\arraystretch}{1.25}
\begin{tabular}{|c|c|c|}
\hline
\(y\) & \(x\) & \(z\) \\ \hline
0 & 462 & \entrybox \\ \hline
\(-3\) & \entrybox & 439 \\ \hline
\(-10\) & 1.25 & \entrybox \\ \hline
3 & \entrybox & 16 \\ \hline
\end{tabular}
\end{center}
}

\newcommand{\qfortyfouranswertable}{%
\begin{center}
\renewcommand{\arraystretch}{1.25}
\begin{tabular}{|c|c|c|}
\hline
\(y\) & \(x\) & \(z\) \\ \hline
0 & 462 & 804 \\ \hline
\(-3\) & 145 & 439 \\ \hline
\(-10\) & 1.25 & 5.5 \\ \hline
3 & 1,679.5 & 16 \\ \hline
\end{tabular}
\end{center}
}

\begin{document}

\packetheading{Lesson 5-4 $\cdot$ Period 3 $\cdot$ Teacher Edition}{Radical Equations --- Assessment-Critical: Solve for Variable + Modeling}

\sectionbanner{OBJECTIVES}
{\small
\renewcommand{\arraystretch}{1.25}
\noindent\begin{tabular}{|p{1.8in}|p{4.8in}|}
\hline
\lessonrow{Math Objective}{Rewrite radical/rational-exponent equations to isolate a specified variable; evaluate such rewritten formulas at given values to produce tables and graphs.}
\lessonrow{Language Objective}{Explain in writing which variable is being isolated and why, using a formula's context (volume, BSA, speed).}
\lessonrow{Essential Understanding}{Solving a formula for a new variable is the SAME algebra as solving a radical equation --- isolate the radical/rational-exponent term, raise to the reciprocal power, clean up. Works for both pure-algebra and real-world cases.}
\end{tabular}
}

\sectionbanner{TOPIC GOALS / ESSENTIAL QUESTION / MATERIALS}
{\small
\renewcommand{\arraystretch}{1.25}
\noindent\begin{tabular}{|p{1.8in}|p{4.8in}|}
\hline
\lessonrow{Topic Goals}{Rewrite formulas to isolate a specified variable (the Topic 5 assessment's core skill); evaluate rewritten formulas at given V/x values.}
\lessonrow{Essential Question}{How do we extract ANY variable from a radical/rational-exponent formula, and why does the same procedure work for both algebra and real-world contexts?}
\lessonrow{Materials}{Student packet, pencil, calculator for table evaluation. Practice \#44 is the direct assessment rehearsal --- treat it like the real test.}
\end{tabular}
}

\begin{callout}{\IconPin\ PRE-FLIGHT --- P3 IS ASSESSMENT-CRITICAL}{calloutpink}
Topic 5 Performance Assessment Items 1A/1B (pyramid volume \(\rightarrow\) solve for x \(\rightarrow\) evaluate) and Item 2 (complete the table) are directly rehearsed by Practice \#10 and Practice \#44.\\
Ex 2 (rewrite a formula --- Golden Gate cables) and q44 (table-completion) are paired in Launch --- students see both `isolate a variable' and `evaluate into a table' back-to-back.\\
No DOK-3 driver today (P2 owned the spine). P3 is pure DOK-2 mastery.\\
Pace: Try-It 2 (5) \(\rightarrow\) \#10 (7) \(\rightarrow\) \#25/\#26/\#27 (5 each) \(\rightarrow\) \#39 (8). Total \(\sim\)35 min.\\
If pace slows: cut \#27 first. Never cut \#10 (Item 1A rehearsal).
\end{callout}

\phasetag{bridge from P2 --- extraneous solutions (HOT)}
\frameworkphaseheader{Do Now}{3}{5}{%
\phaseitems{
  \item Hand out at door.
  \item Circulate 3 min.
  \item Cold-call for extraneous solution reasoning.
}%
}{%
\phaseitems{
  \item Identify the extraneous solution and justify the check step.
  \item One-sentence explanation of why squaring can introduce false solutions.
}%
}{%
\phaseitems{
  \item HOW do you know which solution to reject?
  \item Based on yesterday: predict --- will rewriting a formula for a variable produce extraneous solutions?
}%
}{Solo teacher: circulate silently.}

\bankitem{Do Now (5-4-savvas-q19)}{%
Higher Order Thinking Some, but not all, equations with rational exponents have extraneous solutions. What is the relationship between the exponents and the possibility of having extraneous solutions for equations with rational exponents? Explain your reasoning.
}{}

\teacheranswerbox{\IconCheck\ ANSWER KEY}{%
\textbf{Teacher answer:} Extraneous solutions are introduced when a solving step raises both sides to an even power or undoes an even root, because that can erase sign information. Odd powers and odd roots preserve sign, so they do not create the same false-solution risk. Any candidate solution should still be checked in the original equation.
}

\begin{callout}{\IconSearch\ PREDICTION (spend 30 sec)}{calloutpeach}
Before we rewrite formulas together, predict: if you solved \(V = \frac{1}{3}\pi r^2 h\) for \(r\), would you expect an extraneous solution? Why or why not? (No wrong answers --- just reason it out.) We'll revisit in Launch.
\end{callout}

\sectionbanner{LAUNCH --- Rewrite a Formula + Table Completion (teacher models Ex 2 + \#44)}

\begin{callout}{\IconBook\ ISOLATE-A-VARIABLE RULES}{calloutgreen}
\begin{enumerate}[leftmargin=1.6em,itemsep=2pt,topsep=2pt]
\item Identify which variable to isolate (look at the final answer form).
\item Strip away outer operations LAST-IN, FIRST-OUT (addition/subtraction last \(\rightarrow\) multiplication/division \(\rightarrow\) exponent/radical first to remove).
\item Raise both sides to the RECIPROCAL power to undo a rational exponent (\(m/n \rightarrow n/m\)).
\end{enumerate}
\end{callout}

\begin{callout}{\IconBook\ TABLE-EVALUATION RULES}{calloutyellow}
\begin{enumerate}[leftmargin=1.6em,itemsep=2pt,topsep=2pt]
\item For each input value in the table, substitute into the rewritten formula.
\item Simplify rational exponents step-by-step (use calculator for irrational outputs).
\item Round to specified precision.
\item The table is a point-cloud of the formula's graph --- plot the points to see the shape.
\end{enumerate}
\end{callout}

\phasetag{teacher models Ex 2 (rewrite formula) + \#44 (table) back-to-back}
\frameworkphaseheader{Launch}{2}{12}{%
\phaseitems{
  \item FIRST 6 min: Example 2 --- Golden Gate suspension cable formula. Identify the target variable. Undo operations LAST-IN FIRST-OUT: subtract, then divide, then raise to reciprocal power. Emphasize: \((m/n) \rightarrow (n/m)\) flips the exponent.
  \item NEXT 5 min: Practice \#44 --- table-completion. Show the rewritten formula first, then substitute each V value row-by-row. Demonstrate calculator use for irrational outputs. Emphasize: this is EXACTLY Assessment Item 2.
  \item 30 sec: point at packet rule cards. ``Mark them. These are your assessment reference.''
  \item Do NOT solve Try-It 2 or Explore items. Release at minute 17.
}%
}{%
\phaseitems{
  \item Watch Ex 2. Copy each step. Circle the reciprocal-power move.
  \item Watch \#44. Copy the table. Note which column is input (V) and which is output (x).
  \item Copy BOTH rule cards into margins.
}%
}{%
\phaseitems{
  \item Why do we raise to the RECIPROCAL power to undo \(x^{(2/3)}\)?
  \item In the table: as V increases, what happens to x? Does that make sense for volume?
  \item For Ex 2: what physical constraint rules out negative values of the variable?
  \item How is rewriting a formula different from solving an equation? (It's the same algebra --- just a letter, not a number, on the right side.)
}%
}{Project both examples. Think aloud on the reciprocal-power step.}

\bankitem{Example 2 --- Rewrite a Formula (Golden Gate cables) (5-4-savvas-example-2-lesson-5-4)}{%
\textbf{Rewrite a Formula}

The suspension cables from the Golden Gate Bridge's towers are farther above the roadway near the towers and closer to the roadway near the middle of the bridge. You can figure out your distance from the middle of the bridge, \(x\), in feet, and the height of the suspension cable, \(y\), in feet, by using the given equation. About how far is the cable from the roadway when you are 200 ft from the middle of the bridge?

\visualplaceholder{Golden Gate Bridge with the equation \(x = \frac{\sqrt{y - 220}}{0.010583}\) pointing to the suspension cable.}

Rewrite the equation to isolate \(y\) so you can express the vertical distance between the roadway of the bridge and the suspension cable in terms of \(x\).

\begin{align*}
x &= \frac{\sqrt{y - 220}}{0.010583} \\
0.010583x &= \sqrt{y - 220} \\
(0.010583x)^2 &= (\sqrt{y - 220})^2 \\
0.000112x^2 &\approx y - 220 \\
0.000112x^2 + 220 &\approx y \\
0.000112(200)^2 + 220 &\approx y \\
224.48 &\approx y
\end{align*}

The cable is about 224 ft above the bridge's roadway when you are 200 ft from the middle of the bridge.
}{}

\teacheranswerbox{\IconCheck\ ANSWER KEY}{%
\textbf{Teacher answer:} \(y \approx 0.000112x^2 + 220\), so when \(x = 200\), \(y \approx 224.48\). The cable is about 224 ft above the bridge's roadway.
}

\bankitem{Practice \#44 --- Table Completion \IconStar\ ASSESSMENT REHEARSAL (Item 2) (5-4-savvas-q44)}{%
Complete the table to solve for the unknown value in the equation \(y = \sqrt[3]{2x + z} - 12\), using the given values in each row.

\qfortyfourtable
}{}

\teacheranswerbox{\IconCheck\ ANSWER KEY}{%
\textbf{Answer key (from source):}
\qfortyfouranswertable
}

\begin{callout}{\IconStar\ ASSESSMENT NOTE --- Practice \#44}{calloutyellow}
Practice \#44 is the algebraic twin of Topic 5 Assessment Item 2 --- complete the table for \(x = \sqrt[3]{4V}/2\) at V values \(\{0, 0.25, 2, 16, 32, 64, 128\}\).\\
Any student who can complete \#44 can complete Assessment Item 2.\\
Teacher: walk through at least the first two rows of the table live so students see the calculator workflow for irrational outputs.
\end{callout}

\begin{callout}{\IconMic\ TEACHER SCRIPT}{calloutblue}
1. ``Yesterday we solved radical equations --- one variable, one answer. Today we solve for A variable --- we rewrite the whole formula.''\\
2. Ex 2: ``Target is the length variable. What's in the way? First the `+', then the coefficient, then the exponent. LAST-IN FIRST-OUT. Flip \((2/3) \rightarrow (3/2)\). Done.''\\
3. ``\#44: I already rewrote the formula. Now I just substitute. \(V=0 \rightarrow x=0\). \(V=0.25 \rightarrow x=\sqrt[3]{1}/2 = 0.5\). Calculator for the rest.''\\
4. ``Same procedure, new purpose. The math is identical whether it's a number or a table of numbers on the right.''\\
5. Release.
\end{callout}

\sectionbanner{EXPLORE --- Think-Pair-Share (35 min)}

{\small Work with a partner. Practice \#10 rehearses Assessment Item 1A (isolate x). Practice \#44 (done in Launch) rehearses Assessment Item 2 (complete the table). Plan \(\sim\)8 min for \#39 (BSA application).}

\phasetag{TPS \textperiodcentered\ pure DOK-2 mastery \textperiodcentered\ assessment rehearsal}
\frameworkphaseheader{Explore}{2}{35}{%
\phaseitems{
  \item 6 laps across 35 min.
  \item Laps 1-2: Try-It 2 + \#10. Isolate-a-variable practice. Press reciprocal-power step on every pair.
  \item Laps 3-5: \#25 + \#26 + \#27 --- solve for y with rationalization. Pattern recognition set.
  \item Lap 6: \#39 --- BSA application (modeling). Press interpretation: what does the variable mean?
  \item Do not give answers. Use sentence frame to press interpretation.
}%
}{%
\phaseitems{
  \item 6 items in order. Isolate-a-variable first (Try-It 2, \#10), then rationalization pattern (\#25--\#27), then application (\#39).
  \item For \#10: BEFORE solving, state which variable you're isolating and what power will undo it.
  \item Justify with sentence frame.
}%
}{%
\phaseitems{
  \item What reciprocal power undoes \(x^{(3/2)}\)? What about \(x^{(2/5)}\)?
  \item For \#25--\#27: what's the same in every problem? What changes?
  \item For \#10: how does your answer connect to Assessment Item 1A?
  \item For \#39: what do the units of BSA tell you about whether your answer is reasonable?
  \item After \#10: could you now complete a table like \#44 using your answer? Try one row.
}%
}{Solo teacher: 6 laps. Press the reciprocal-power step on every item. \#10 IS the Item 1A rehearsal.}

\textbf{Bank-item labels (in order):}
\begin{enumerate}[leftmargin=1.6em,itemsep=2pt,topsep=2pt]
  \item Try It 2
  \item Practice \#10 \IconStar\ ASSESSMENT ITEM 1A
  \item Practice \#25
  \item Practice \#26
  \item Practice \#27
  \item Practice \#39 --- BSA Application
\end{enumerate}

\bankitem{Try It 2 (5-4-savvas-try-it-2-lesson-5-4)}{%
The speed, \(v\), of a vehicle in relation to its stopping distance, \(d\), is represented by the equation \(v = 3.57\sqrt{d}\). What is the equation for the stopping distance in terms of the vehicle's speed?
}{}

\teacheranswerbox{\IconCheck\ ANSWER / ANTICIPATED TALK}{%
\textbf{Answer key (from source):} \(d = \left(\frac{v}{3.57}\right)^2\)
}

\bankitem{Practice \#10 \IconStar\ ASSESSMENT ITEM 1A (5-4-savvas-q10)}{%
Rewrite the equation \(y = \sqrt{\frac{x - 48}{6}}\) to isolate \(x\).
}{}

\teacheranswerbox{\IconCheck\ ANSWER / ANTICIPATED TALK}{%
\textbf{Teacher answer:} \(x = 6y^2 + 48\).
}

\bankitem{Practice \#25 (5-4-savvas-q25)}{%
Solve for \(y\). \(x = 3\sqrt[3]{15 + y}\)
}{}

\teacheranswerbox{\IconCheck\ ANSWER / ANTICIPATED TALK}{%
\textbf{Answer key (from source):} \(y = \frac{x^3}{27} - 15\)
}

\bankitem{Practice \#26 (5-4-savvas-q26)}{%
Solve for \(y\). \(x = \frac{\sqrt{2y}}{26}\)
}{}

\teacheranswerbox{\IconCheck\ ANSWER / ANTICIPATED TALK}{%
\textbf{Answer key (from source):} \(y = 338x^2\)
}

\bankitem{Practice \#27 (5-4-savvas-q27)}{%
Solve for \(y\). \(x = \frac{\sqrt{y - 14.2}}{0.05}\)
}{}

\teacheranswerbox{\IconCheck\ ANSWER / ANTICIPATED TALK}{%
\textbf{Answer key (from source):} \(y = 0.0025x^2 + 14.2\)
}

\bankitem{Practice \#39 --- BSA Application (5-4-savvas-q39)}{%
Solve using the formula \(BSA = \sqrt{\frac{H \cdot M}{3{,}600}}\). A sports medicine specialist determines that a hot-weather training strategy is appropriate for this athlete. To the nearest tenth, what can the mass of the athlete be for the training strategy to be appropriate?

\visualplaceholder{[IMAGE: Athlete in white coat at a finish line with labels ``BSA < 2.0'' and ``165 cm.'']}
}{}

\teacheranswerbox{\IconCheck\ ANSWER / ANTICIPATED TALK}{%
\textbf{Answer key (from source):} less than \(87.3\) kg
}

\sectionbanner{SHARE / SUMMARY (5 min)}

\phasetag{Concept Summary + Topic 5 Performance Task preview}
\frameworkphaseheader{Share / Summary}{2}{5}{%
\phaseitems{
  \item Cold-call 1 pair to present \#10 (isolate x) using the sentence frame.
  \item Project Concept Summary (from student packet).
  \item ``Topic 5 Performance Assessment: Items 1A/1B ask you to rewrite a pyramid-volume formula for x, then evaluate. Item 2 asks you to complete a table. You just rehearsed BOTH.''
}%
}{%
\phaseitems{
  \item Presenting pair uses sentence frame for \#10.
  \item Rest of class signals and annotates their Concept Summary.
}%
}{%
\phaseitems{
  \item What are the TWO big ideas from Lesson 5-4 P3?
  \item If you see a rational exponent on the assessment, what's the FIRST thing you do?
}%
}{Frame today as assessment rehearsal. Students leave confident on isolating variables.}

\begin{sentenceframebox}
\textbf{Sentence frames:}

Pair share (\#10 or \#39): ``To isolate \_\_\_, I first \_\_\_, then raise both sides to the power \_\_\_. The rewritten formula is \_\_\_. Plugging in \_\_\_ = \_\_\_ gives \_\_\_. In context this means \_\_\_.'' 

Whole class: one pair reads their answer. Class signals thumbs up/sideways.
\end{sentenceframebox}

\begin{callout}{\IconChart\ CONCEPT SUMMARY (your reference for Topic 5 assessment)}{calloutell}
ISOLATING A VARIABLE FROM A RADICAL/RATIONAL-EXPONENT FORMULA:\\
\hspace*{1.5em}--- Identify the variable to isolate and the form of the final answer.\\
\hspace*{1.5em}--- Undo outer operations first (add/sub), then multiply/divide, then remove the exponent/radical.\\
\hspace*{1.5em}--- To undo \(x^{(m/n)}\): raise both sides to the power \((n/m)\).\\
\hspace*{1.5em}--- Simplify: distribute, combine like terms, solve for the target variable.

\medskip
COMPLETING A TABLE FROM A REWRITTEN FORMULA:\\
\hspace*{1.5em}--- Use your rewritten formula (e.g., \(x = \sqrt[3]{4V}/2\)) as the rule for each row.\\
\hspace*{1.5em}--- Substitute each V-value. Use a calculator for irrational outputs. Round as directed.\\
\hspace*{1.5em}--- Check: does the table's pattern match the shape you expect (growth/decay/root curve)?\\
\hspace*{1.5em}--- Assessment Item 2 expects THIS process --- from the rewritten formula to the table.
\end{callout}

\sectionbanner{EXIT TICKET}

\phasetag{pre-assessment confidence check}
\frameworkphaseheader{Exit Ticket}{1-2}{3}{%
\phaseitems{
  \item Collect at door.
  \item Tonight: scan for students who can't complete the reciprocal-power step or the table row. Insert a 5-min Do Now recap when the Performance Assessment lands.
}%
}{%
\phaseitems{
  \item Complete the three fill-in stems.
  \item Write one biggest thing learned.
}%
}{%
\phaseitems{
  \item Did any student leave the reciprocal power blank?
  \item Did any student evaluate the formula correctly but skip the table interpretation?
}%
}{Collect. No grade. Use as data for assessment prep.}

\begin{callout}{\IconClipboard\ EXIT TICKET --- Student Form}{calloutyellow}
To isolate a variable from a rational-exponent formula I \_\_\_.\\
When I evaluate the formula at a specific x, I \_\_\_.\\
One biggest thing I learned today: \_\_\_\_\_\_\_\_\_\_\_\_\_\_\_\_\_\_\_\_\_\_
\end{callout}

\clearpage

\begin{callout}{\IconClock\ PERIOD A / PERIOD F --- 65-MIN CLASS NOTES}{calloutpink}
Both sections should have 65 min. Use extra 10 min on \#10 and \#39 (assessment-critical).\\
If finished early: Practice \#46 (Escape Velocity) as fast-finisher stretch.\\
Alternative: project \#44 as a second table-completion prompt --- more assessment rehearsal.
\end{callout}

\begin{callout}{\IconPuzzle\ MODIFICATIONS / SUPPORT FOR IEP STUDENTS}{calloutiep}
\textbullet\ Provide both Rule cards as half-sheet cut-outs on their desks.\\
\textbullet\ For \#10, provide a scaffold showing the first two steps already completed so they focus on the reciprocal-power move.\\
\textbullet\ Accept verbal formula-rewriting in lieu of written for \#10.
\end{callout}

\begin{callout}{\IconGlobe\ SUPPORTS FOR ELL STUDENTS}{calloutell}
\textbullet\ Sentence frames on student page (don't skip).\\
\textbullet\ Word cards: ``isolate'' = get the variable alone; ``reciprocal power'' = flip the fraction in the exponent.\\
\textbullet\ ``rewrite'' = change what the formula solves for, without changing the math.\\
\textbullet\ ``table'' = list of (input, output) pairs --- each row is one substitution.\\
\textbullet\ Pair ELL students with English-dominant partners for TPS.
\end{callout}

\end{document}
